% Level 2, Week 10 homework -- project work
\homeworkweek{10}{Project Studio 1}{15--20 minutes. Do Core first; Challenge and Creative are optional.}
\namedateline

\begin{tip}[No computer at home this week?]
That is fine. Do the paper parts (the dev log, the Challenge plan and the Creative sketch) and bring
them next week. Your project files stay in your team folder; never email or upload them.
\end{tip}

\begin{core}[Finish your own sticky notes (about 10 minutes)]
\begin{enumerate}
  \item Open your team's project in Thonny. Pick \textbf{one} sticky note from your TO DO column that is yours.
  \item Build it in small steps. Press \keys{F5} after each step. Save with \keys{Cmd}+\keys{S} often.
  \item If you get a red error, read the last line of the traceback and the line number it points to. Stuck for 5 minutes? Write down what you tried and stop: that is useful too.
  \item Write a short \textbf{dev log} (developer diary) below.
\end{enumerate}
What I worked on: \answerlines{1}
What broke, and how I fixed it (or what I tried): \answerlines{1}
What I will do next: \answerlines{1}
\end{core}

\begin{challenge}[Tests on paper, before the code]
Next week you will write pytest tests. Get ahead: choose one of your functions that \texttt{return}s
something and write three examples of what goes in and what should come out. Make at least one an
\textbf{edge case} (a value at the very limit, like the edge of the screen or a score of 0).

\renewcommand{\arraystretch}{1.6}
\begin{tabularx}{\linewidth}{|X|X|X|}
\hline
\textbf{Function and input} & \textbf{Expected answer} & \textbf{Normal or edge?} \\ \hline
\texttt{next\_x(0, 1)} & \texttt{20} & normal \\ \hline
 & & \\ \hline
 & & \\ \hline
 & & \\ \hline
\end{tabularx}
\end{challenge}

\begin{creative}[Title screen sketch]
Design the title screen for your project on squared paper: its name, a picture, and one line that
explains the keys (\enquote{← → to move, space to start}). Use colours from the Okabe--Ito palette
(orange, sky blue, bluish green, yellow, blue, vermillion, reddish purple) and make sure nothing
depends on colour alone. Use your nickname, not your real name.
\answerspace{5cm}
\end{creative}

\begin{selfcheck}
  \item I worked on one of my own tasks and saved my work.
  \item My program still runs (or I wrote down exactly where it breaks).
  \item I wrote my dev log: what I did, what broke, what is next.
  \item Anything I copied or remixed has a credit comment.
\end{selfcheck}
